\chapter*{Preface}
\addcontentsline{toc}{chapter}{Preface}
\markboth{Preface}{Preface}

This book is about a single, narrow claim: that the information needed to
expose a C++ library to Python, to JavaScript, to Lua, to a REST endpoint or
to a property panel is \emph{already in the header}, and that a compiler can
be made to read it out.

Everything else follows from taking that claim seriously. You do not write a
binding. You write down \emph{what} you want bound --- a list of headers and a
list of target languages --- and a program compiled against your types answers
the rest by reflection. Your class definitions never change.

\section*{Who this is for}

You have a C++ library. You want it callable from somewhere else, ideally from
several somewheres, and you would rather not maintain a hand-written wrapper
per destination. You are comfortable with CMake and with your own build.

You do \emph{not} need to know anything about C++26 static reflection to use
rosetta. Reflection is what makes the tool possible, and Chapter~\ref{ch:pipeline}
explains where it runs and why the pipeline has the shape it does --- but the
surface you work with is a JSON file and a command-line tool.

\section*{How to read it}

\begin{description}[leftmargin=1.4cm,style=nextline,font=\normalfont\itshape]
  \item[In a hurry] Chapter~\ref{ch:first} is a complete first binding, from an
    unmodified header to an importable Python module. Read that, then come back.
  \item[Evaluating rosetta] Chapters~\ref{ch:what} and~\ref{ch:pipeline} say what
    the system does and what it costs, including the parts that do not work.
  \item[Using it in earnest] Part~II is the manifest, field by field, and is the
    part you will return to. Part~IV is the tool that consumes it.
  \item[Documenting what you bind] Part~III is the two ways metadata reaches a
    binding: annotations you write, and the Doxygen comments and default
    arguments your headers already carry
    (chapter~\ref{ch:documentation}, on by default).
  \item[Pushing past bindings] Part~V covers the runtime object model, scripting
    the metadata itself, and writing a backend of your own.
\end{description}

The appendices are lookup tables: every manifest field, every annotation, every
target, and the errors you are most likely to hit.

\section*{Conventions}

\code{Monospace} is a literal --- a field name, a path, a command. A
\emph{manifest} always means the \code{manifest.json} driving one binding
project. \emph{Generation host} is the machine that runs the reflection
compiler; \emph{target} is wherever the generated binding is finally built,
which is deliberately not the same machine.

Two callouts appear throughout:

\begin{note}
Something worth knowing that is not on the shortest path.
\end{note}

\begin{warn}
Something that will cost you an afternoon if you meet it unprepared.
\end{warn}

\section*{Status}

Rosetta is a prototype, and this book says so where it matters. It tracks
in-flight C++ proposals --- P2996 (reflection), P3394 (annotations), P3294
(token injection) --- and is developed against the Bloomberg
\code{clang-p2996} fork, which is currently the only front-end that carries
all of what it uses. What the tool \emph{generates} has no such constraint:
that is the whole point, and Chapter~\ref{ch:pipeline} is where the distinction
is made precise.

\section*{Related documents}

The repository carries the reference material this book is built from ---
\code{docs/MANIFEST.md}, \code{docs/ROSETTA\_GEN.md},
\code{docs/ANNOTATIONS.md}, \code{docs/COVERAGE.md} and the rest. Where a
detail here is compressed, those are the long form. The academic write-up of
the same system, with its evaluation and its related work, lives under
\code{paper/}; it argues that reflection can serve as an interface description
language, where this book simply assumes it and gets on with the job.
